\documentclass{article}
\usepackage[french]{babel}
\usepackage[T1]{fontenc}
\usepackage[utf8]{inputenc}
\usepackage{amsmath}

\title{Analyse spectrale des op\'erateurs diff\'erentiels}
\author{M.~Dupont}
\date{2026}

\begin{document}
\maketitle

\section{Introduction}

L'analyse spectrale des op\'erateurs diff\'erentiels est un domaine fondamental
des math\'ematiques appliqu\'ees.  Dans cet article, nous \'etudions les propri\'et\'es
de convergence des valeurs propres pour une classe d'op\'erateurs elliptiques.

% LANG-003 / FR-007: missing NBSP before ; : ! ?
Le th\'eor\`eme principal est le suivant: pour tout op\'erateur $A$ born\'e, le
spectre est compact. C'est un r\'esultat important; il g\'en\'eralise les travaux
ant\'erieurs! Peut-on l'am\'eliorer?

% Correct: with NBSP before punctuation
Le r\'esultat est remarquable~; il unifie plusieurs approches~! La question est~:
comment l'appliquer en pratique~?

% LANG-005: wrong quote style (straight quotes instead of guillemets)
La m\'ethode dite "de Galerkin" est bien connue.  On parle aussi de "convergence
faible" dans ce contexte.

% Correct: guillemets
La m\'ethode dite \og de Galerkin\fg{} est classique.  On parle de \og convergence
forte\fg{} lorsque la norme est pr\'eserv\'ee.

% FR-007: Euro sign without thin space
Le co\^ut du calcul est estim\'e \`a 50\euro{} par heure de processeur.
L'abonnement revient \`a 120\euro{} par mois.

% Correct: with thin space
Le budget total est de 500\,\euro{} pour le projet.

% LANG-011: oe vs \oe{}
Le coeur du probl\`eme r\'eside dans l'approximation.  L'oeuvre de Lebesgue
est fondamentale pour cette th\'eorie.

% Correct: \oe ligature
Le c\oe{}ur de l'algorithme repose sur la d\'ecomposition.  L'\oe{}uvre de
Cauchy inspire notre approche.

\section{R\'esultats}

Soit $T \colon H \to H$ un op\'erateur lin\'eaire born\'e sur un espace de Hilbert.
On d\'emontre que $\sigma(T) \subset \overline{B(0,\|T\|)}$.

\begin{theorem}
Pour tout $\epsilon > 0$, il existe $N \in \mathbb{N}$ tel que pour tout $n \ge N$,
on a $|\lambda_n - \lambda| < \epsilon$.
\end{theorem}

La d\'emonstration utilise le th\'eor\`eme de Banach-Steinhaus et le lemme de Riesz.

\section{Conclusion}

Les r\'esultats obtenus permettent de caract\'eriser compl\`etement le spectre des
op\'erateurs consid\'er\'es.  Les applications \`a la m\'ecanique quantique et \`a
la th\'eorie du contr\^ole sont imm\'ediates.

\end{document}
